\chapter{Ordinary actual Alice angle, with arbitrary state ratio}
\label{component:p9f00f949fc8c}
\noindent\textbf{Component }\path{p9f00f949fc8c}\par
{\small\noindent Source: \path{R13/baseline/baseline/baseline/baseline/baseline/baseline/advances/c1/arbitrary_state_ratio.tex}\par}
{\small\noindent SHA256: \path{9f00f949fc8c2e401f57d03bb695867f73c6581778eb8351d4f5dcbdfa388eba}\par}
{\small\noindent Inventory use: entire scoped proof | C1\_\allowbreak{}BARE\_\allowbreak{}OTHER\_\allowbreak{}TERMINALS; precise subordinate uses retained below\par}
\medskip
\section{Ordinary actual Alice angle, with arbitrary state ratio}
\label{p9f00f949fc8c:sec:c1-arbitrary-state-ratio}

The positive bounded-unit assumption on $s=g/f$ can be removed from
the arbitrary-opening C1 theorem. The new ingredient is a different
grading of the already proved polynomial contact--conductor identities.
No new ideal membership or division is required.

\begin{theorem}
Use the unchanged native C1 source, normalized by $c=2$ and
$\lambda=n=0$. Suppose $k$ is a positive bounded unit, $f,g>0$ are
infinitesimal, and
\[
 \begin{gathered}
 \eta=f^2,\quad\zeta=g^2,\quad K=1+k^2,\qquad
 R=K-ak^2\eta,\quad T=K-k^2t\zeta,\\
 a,t>0,\qquad M=R-t>0,\qquad J=T-a>0.
 \end{gathered}
 \tag{AS1}
\]
Then the full original parent is tensor SOS at every independent
projective test and every real or complex state. The ratio $s=g/f$
is arbitrary positive. No positive lower bound on $a,t,M,J$ and no
boundedness of original event coefficients is assumed.
\end{theorem}

The native source and every retained packet are exactly those of
(OA2), (AO3)--(AO5), with $s=g/f$ and $\Delta=s/k$. Since
$a<T<K$ and $t<R<K$, the shapes are bounded, while $R,T$ are units
of residue $K$. The six coquotients $c_2,c_b,c_3,c_5,c_4,c_\alpha$
have residue $2K$, and $c_d$ has residue $4K^2$.
All original events are positive. For the two subtracted expressions,
this follows from
\[
 \frac Da=\frac1J\left(\frac{2c_5}{c_3}
                   -J\eta\frac{c_d}{c_3}\right),\qquad
 \frac\beta a=\frac1{s^2t}
      \left(\frac{2k^2c_b}{c_3}-t\zeta\frac{c_d}{c_3}\right).
 \tag{AS2}
\]
The parentheses have positive unit residues $2$ and $2k^2$.

\paragraph{Exact reflection reduces to bounded $s$.}
The whole-parent reflection is
\[
 u'=1/u,\quad v'=v,\quad
 (W',X',Y',Z')=(vW-X,-W-vX,Z-vY,Y+vZ),
 \tag{AS3}
\]
with $(\lambda\ n)(m\ c)(A\ \alpha)(D\ \beta)$ on events and
\[
 u^2F(\mathrm{ref}(c);1/u,v;\psi')=(1+v^2)F(c;u,v;\psi).
 \tag{AS4}
\]
This is an identity of the full parent, including the signed core.
At the native contact set
$H_f=1+(1-a)\eta$, $H_g=1+(1-t)\zeta$, and
$V=1+k^2\eta\zeta$. Both $H_f,H_g$ are positive units. After
normalizing the reflected contact by its positive $Z'$ coordinate,
the exact reflected native parameters are
\[
\begin{gathered}
 k'=1/k,\qquad s'=\frac{H_g}{sH_f},\qquad
 \eta'=\frac{k^2\zeta H_f}{H_g},\qquad
 \zeta'=\frac{k^2\eta H_g}{H_f},\\
 a'=\frac{M}{k^2H_f},\quad t'=\frac{J}{k^2H_g},\qquad
 M'=\frac{aV}{k^2H_g},\quad J'=\frac{tV}{k^2H_f}.
\end{gathered}\tag{AS5}
\]
For example, the reflected state has
$r'=Q_c/H_g$, $\theta'=\eta J/H_f$, where
$Q_c=1-k^2\zeta(1-a\eta)$; the differences
$H_g-Q_c=\zeta M$ and $H_f-E=\eta J$ prove (AS5).
The reflected coefficient vector, after positive normalization to
$c'=2$, is the reflected native source: this follows from the same
six contact equations and their nonzero selected minor used for the
original native source. All displayed shapes and gaps are positive.
Thus unbounded $s$ becomes infinitesimal $s'$, with both
$\eta',\zeta'$ infinitesimal and $k'$ a unit. Unit $s$ is the
previous arbitrary-opening theorem. It remains to prove the case
$s$ infinitesimal.

\paragraph{Outer angles require only a positive current mass.}
Put $D_1=2c_5/c_3-J\eta c_d/c_3$, a unit of residue $2$.
For $g_c=cD/(c+D)$ the exact cancellation
\[
 \frac{g_c}{a}=\frac{2D_1}{2J+aD_1},\qquad
 2J+aD_1=2T+a(D_1-2)
 \tag{AS6}
\]
gives
\[
 d/c,\ d/D\text{ infinitesimal},\qquad
 \gamma=\frac{Ag_c}{d^2}\text{ a positive bounded unit},
 \qquad\bar\gamma=\frac1{4\bar K^3}.
 \tag{AS7}
\]
The two-strong-current identity therefore covers every test with
$1/v$ bounded, including Bob infinity, for every Alice angle and state.
For $v$ infinitesimal and Alice unbounded, use the exact critical-current
identity with residue diagonal
\[
 \operatorname{diag}(1,2,2,4\bar\gamma-\bar c_u^2),
 \qquad c_u=(1+u^2)^{-1/2}.
 \tag{AS8}
\]
Here $\bar c_u=0$ and $\bar\gamma>0$, so it is strict.
The former sufficient assumption $\bar\gamma\ge1/4$ is unnecessary
on this particular reciprocal-Alice chart. Its same exact identity
and homogeneous normalization also include literal Alice infinity.
Consequently only bounded $u$ and infinitesimal $v$ remain.

\paragraph{The complementary state chart.}
Normalize the testing state by a coordinate of maximal absolute value.
Put $D_u=1+u^2$ and $w=(W+uY)/D_u$, a pure testing coordinate.
If $w$ is a unit, the retained-current identity (UW2)--(UW3) applies
with $Q=W+uY=D_uw$ a unit. It requires only bounded $u,v,Z$, not
infinitesimal $u$ or unit $k$. The five roots are
\[
 r_2,\ B,\ \Delta r,\ \Delta r_5,\ \Delta\eta Y.
\]
Their positive diagonal coefficients are
$H_2,H_bD_v,H_3D_v,H_5,4RTc_4D_uD_v$, all units.
The only two negative corrections have factors
$\Delta Z/Q$ and $\Delta^2\eta Z^2/Q^2$, hence zero residue.
Restoring the retained packets proves this chart.

Otherwise $w$ is infinitesimal. Literal axes and a nonnegative pure
testing core are direct original-event SOS exits. On the negative-core
branch, sign gauges give $u,v>0$. The coordinate $Z$ must be a unit:
at $v$ of zero residue the
core has residue $2((X+uZ)^2+(uW-Y)^2)$. If $Z,w$ and these two
currents all had zero residue, the normalized state would have every
coordinate of zero residue. Thus a negative-core test with small $w$
has unit $Z$; normalize $Z=1$ by a pure testing factor. The
pure-testing Schur identity now gives $w>0$; before this normalization
its conclusion is $wZ>0$.

\paragraph{The same conductor, with $k$ of weight zero.}
In this last chart use exactly (SD3)--(SD8), retaining all four
positive packets. Set
\[
 \delta=\Delta R/K,\qquad b^2=\eta K/R,
 \qquad Q=K+ak^2b^2,\quad T_n=K^2-t\delta^2k^4Qb^2,
 \quad\mathcal C=K^6Q^4T_n.
\]
Here $\delta,b$ are positive infinitesimals, and $k,a,t,u$ are bounded.
The exact source substitution is unchanged:
\[
 \eta=Kb^2/Q,\quad\Delta=\delta Q/K,\quad
 \zeta=\delta^2k^2Qb^2/K,\quad R=K^2/Q,\quad T=T_n/K.
\]
Let $E_0,L_1,L_2$ be the scalar correction and two full first normals
in the contact--conductor proof. Its polynomial construction places
\[
 E_0\in(x,y,z,h,q_1)^2,\qquad L_1,L_2\in(x,y,z),
 \tag{AS9}
\]
where
\[
 \begin{gathered}
 A_0=wv,\ B_0=\delta v,\ C_0=\delta b^2uw,\ D_0=\delta^2k^3b^2,\\
 x=B_0-C_0,\quad y=A_0-kC_0,\quad z=uB_0-kC_0,\quad h=kC_0-D_0.
 \end{gathered}
\]
These are universal polynomial identities; their proof never assumed
that $k$ was small or inverted it.

Give $(\delta,k,b,a,t,u,w,v)$ weights $(1,0,1,0,0,0,1,3)$.
All four monomials $A_0,B_0,C_0,D_0$, and all the soft roots, have
weight $4$. The substitutions at contact
$(u,w,v)=(k,\delta k,\delta k^2b^2)$ preserve this grading.
The factored five-term scalar gives
\[
 \operatorname{wt}E_0\ge10,\qquad
 \operatorname{wt}L_1\ge8,\qquad
 \operatorname{wt}L_2\ge6.\tag{AS10}
\]
Here is a short way to read these bounds. One has
$Q-K=O(2)$, $T_n-K^2=O(4)$, $V_n-K^2=O(4)$,
$D_v-1=O(6)$, $Y_n=O(1)$ and $B_n+KT_n=O(4)$.
The six cleared coquotients equal $2K^3Q^3+O(2)$, while
$C_d=4K^4Q^3+O(2)$. In the five-term scalar formula, discard terms
of weight at least $10$; its four weight-eight squares and cross term
are precisely $8K^2D_u\mathcal C P$. For the normals use
\[
 K(C_b-C_2)=\delta^2k^2Qb^2C_3,
 \qquad Q(C_3-C_5)=-k^2Kb^2C_b,
\]
and $S_B,S_C=O(2)$ in their displayed factorizations. This gives
(AS10) without an expanded source calculation.

Monic contact division preserves these weights. The conductor lifts
therefore have coefficient weights at least $2,4,2$, respectively.
Every such polynomial coefficient contains a positive power of
$\delta,b,w$, or $v$, and has zero residue. There are no negative
exponents, including on $a,t$; the only source denominators used before
this polynomial construction are the displayed positive bounded units.

The leading scalar remains exactly
\[
 P=D_u(A_0^2+B_0^2)+K(C_0^2+D_0^2)-2K^2B_0C_0
   =x^2+g_4^2+\sum_{i=1}^4a_i^4-4\prod_{i=1}^4a_i.
\]
Here $g_4=uwv-\delta^2k^4b^2$ and
\[
 \begin{gathered}
 a_1=\sqrt{wv},\quad a_2=\sqrt{\delta uv},\quad
 a_3=b\sqrt{\delta kuw},\quad a_4=\delta k^{3/2}b,
 \\
 q_1=a_1a_2-a_3a_4,\qquad q_2=a_1a_3-a_2a_4.
 \end{gathered}
\]
Their radicals separate between coefficient and testing fields;
$y=a_1^2-a_3^2$, $z=a_2^2-a_3^2$, $h=a_3^2-a_4^2$.
The universal two-pairing circuit identity gives a strict constant Gram
on $(x,y,z,h,q_1,q_2)$. Retain $g_4^2$ as a positive square; its
coefficients need not have zero residue when $k,u$ are units.
The complete hard block on $(N,\delta M_B)$ has positive residue
$16\bar K^{14}I_2$ after multiplication by $\mathcal C$.
Indeed $\Delta/\delta=Q/K\sim1$, $H_i\sim8K^2$, and the remaining
hard-block terms have factors $\eta$ or $v$. The full first normals
are exactly those in (AS9), not only their restrictions at contact.
Thus (AS9)--(AS10) give a bounded strict eight-root remainder Gram.
Restore $g_4^2$ and the four positive packets, then undo the positive
pure source multiplier $\mathcal C\Delta^2\Lambda$.

\paragraph{Exhaustion and precise advance.}
The source reflection, unit-$s$ theorem, two outer-angle identities,
unit-$w$ current identity, and small-$w$ conductor cover every source
and every independent projective test under (AS1). Pure testing
normalizations and pure source multipliers are restored separately;
no mixed inverse or radical occurs. Complex states split into real
and imaginary parts.

Together with the complete infinitesimal-$k$ theorem and the exact
orientation $0<k\le1$, this proves the all-tests assertion whenever
the actual native entry has both $\eta$ and $\zeta$ infinitesimal,
for every positive $k$ in that orientation and every state ratio.
The remaining native complement is therefore: $k$ infinitesimal with
$\zeta$ infinitesimal and $\eta$ a unit or unbounded; $k$
infinitesimal with $\zeta$ non-infinitesimal; and $k$ a unit with
at least one of $\eta,\zeta$ non-infinitesimal. Each is still
intersected with physicality and the complements of all other inherited
admissions and symmetry images. No whole-global completion is asserted.
